\section{Optimal-Space Vector Representation}
\label{sec:vector-representation}

\subsection{The Representation Problem}

Consider a vector
\[
    A[1\ldots n],
    \qquad
    A[i]\in\Sigma,
\]
where the alphabet has size \(\Sigma\). The vector contains
\(n\log_2\Sigma\) bits of information, so an injective binary
representation needs at least
\[
    \left\lceil n\log_2\Sigma\right\rceil
\]
bits. Treating the entire vector as one integer in \([\Sigma^n]\) attains
this bound, but reading or changing one coordinate then requires a global
base conversion. At the other extreme, storing every coordinate in
\(\lceil\log_2\Sigma\rceil\) bits gives local access but incurs linear
redundancy whenever \(\Sigma\) is not a power of two.

The objective is to attain the exact information-theoretic bound while
supporting both reads and writes in constant time. The information carrier
lemma from Section~\ref{sec:information-carriers} supplies the required local
primitive: it commits an old spill to binary memory with low redundancy while
leaving a bounded new spill.

\subsection{Grouping Symbols into Carriers}

Partition consecutive coordinates into groups and interpret each group as a
mixed-radix value
\[
    x_i\in[X].
\]
The paper chooses
\[
    X=\Theta(n^2).
\]
Thus a carrier contains \(\Theta(\log_{\Sigma}n)\) original symbols, and the
number of carriers is
\[
    N=O\left(\frac{n}{\log_{\Sigma}n}\right).
\]
Boundary groups are assigned their deterministically known ranges; this does
not change the asymptotic calculation.

This grouping is what amortizes the rounding loss. One application of the
information carrier lemma has redundancy
\[
    O\left(\frac{1}{\sqrt X}\right).
\]
For \(X=\Theta(n^2)\), this is \(O(1/n)\) per carrier, and hence
\[
    N\cdot O\left(\frac{1}{n}\right)=o(1).
\]
Rather than rounding each small alphabet symbol separately, the construction
performs one rounding operation on a carrier containing logarithmically many
symbols, making the sum of all rounding losses vanish asymptotically.

\subsection{A Sequential First Attempt}

Let the carriers be \(x_1,x_2,\ldots,x_N\). Starting with a null spill
\(y_0\), represented by a singleton range, step \(i\) applies the information
carrier mapping
\[
    (x_i,y_{i-1})\longmapsto(m_i,y_i).
\]
Here \(m_i\) is written to an \(M_i\)-bit memory field, and \(y_i\in[Y_i]\)
is passed to the next step. The lemma ensures \(Y_i=O(\sqrt X)\), even though
its exact value may depend irregularly on \(Y_{i-1}\). After the last step,
the final spill \(y_N\) is stored separately.

\subsection{Redundancy of the Sequential Construction}

The difference between the stored length and the information in the carriers
is
\[
    \left\lceil\log_2Y_N\right\rceil
    +\sum_{i=1}^{N}M_i
    -N\log_2X.
\]
Using the chain rule for redundancy, this is at most
\[
    1+
    \sum_{i=1}^{N}
    \left(
        M_i+\log_2Y_i-\log_2Y_{i-1}-\log_2X
    \right).
\]
Each parenthesized term is precisely the redundancy of one information
carrier transformation. The positive \(\log_2Y_i\) term at one step cancels
the negative term at the next, while the leading one accounts for rounding
up the space used by the final spill. Consequently the total redundancy is
\[
    1+N\cdot O\left(\frac{1}{\sqrt X}\right)=1+o(1).
\]
Crucially, intermediate spills are passed forward rather than independently
stored in rounded-up binary fields. They therefore do not contribute one
rounding bit per step.

\subsection{Constant-Time Reads and Writes}

Suppose a requested coordinate lies in \(x_j\). If \(j<N\), the asymmetric
decoding guarantee of the lemma recovers the old spill \(y_j\) from
\(m_{j+1}\) alone. The pair \((m_j,y_j)\) then recovers \(x_j\), from which
the coordinate is extracted. For \(j=N\), the separately stored final spill
\(y_N\) replaces the nonexistent \(m_{N+1}\). Thus decoding never begins at
\(y_0\) and touches only a constant number of memory fields.

To change \(x_j\), first re-encode step \(j\). Its new output spill \(y_j\)
may differ, so step \(j+1\) is also re-encoded when \(j<N\). In the carrier
mapping of Section~\ref{sec:information-carriers}, the new spill is the
quotient part of the carrier and is independent of the old spill. Since
\(x_{j+1}\) has not changed, this second encoding can retain the same
\(y_{j+1}\); only its memory content changes. Propagation therefore stops
after one neighboring step. When \(j=N\), only \(m_N\) and the separately
stored \(y_N\) are changed. Both reads and writes take \(O(1)\) time.

\subsection{Why the Sequential Construction Is Insufficient}

The sequential construction is correct, space-efficient, and locally
accessible. Its defect is instead the amount of nonuniform precomputation.
The spill universes \(Y_i\) may all differ, so different steps may require
different values of \(M_i\), spill-size parameters, information carrier
constants, memory offsets, and fixed division or multiplication constants.
It may therefore require
\[
    N=O\left(\frac{n}{\log n}\right)
\]
sets of precomputed parameters, rather than the \(O(\log n)\) word constants
required by the theorem.

\subsection{The Implicit Binary-Tree Representation}

The final construction places the carriers in an implicit binary tree using
heap numbering:
\[
    \operatorname{left}(i)=2i,
    \qquad
    \operatorname{right}(i)=2i+1,
    \qquad
    \operatorname{parent}(i)=\lfloor i/2\rfloor.
\]
No pointers are stored. A node combines the spills received from its two
children into a mixed-radix value \(y_i\), and uses its own carrier in the
mapping
\[
    (x_i,y_i)\longmapsto(m_i,s_i).
\]
The memory content \(m_i\) is written in the node's assigned memory area, and
the new spill \(s_i\) is passed to its parent. A missing child contributes the
unique value in a singleton, deterministically known range; leaves therefore
have a fixed null input range. This handles incomplete boundary subtrees
without pointers or an auxiliary data structure. The root spill \(s_1\) is
stored separately.

\subsection{Queries and Updates in the Tree}

For a nonroot node \(v\), the parent's memory content independently reveals
the combined old spill at the parent. Its mixed-radix components include the
spill \(s_v\) sent by \(v\). Combining \(s_v\) with \(m_v\) recovers \(x_v\),
and the desired coordinate is then extracted. For the root, the separately
stored \(s_1\) and \(m_1\) recover \(x_1\). Although spills move upward while
the representation is encoded, a query never follows a root-to-node path and
does not recurse; it examines only the node and its parent.

An update re-encodes the target node. Its outgoing spill may change, so the
parent is re-encoded as well. The parent's carrier remains fixed, and the
information carrier mapping consequently retains the parent's outgoing
spill while absorbing the changed child spill into new memory content.
Propagation stops there. Updating the root changes only its memory content
and separately stored spill. Hence both operations take \(O(1)\) time.

\subsection{Reducing the Number of Precomputed Constants}

Parameter sharing is the reason for using the tree. At a fixed level, for an
appropriate height \(k\), node subtrees occur in only three classes:
\begin{enumerate}[label=(\arabic*)]
    \item a prefix of nodes has complete binary subtrees of height \(k\);
    \item at most one node has an incomplete subtree;
    \item the remaining nodes have complete binary subtrees of height
          \(k-1\).
\end{enumerate}
Nodes in the same class and level have the same subtree size. By induction,
their two child spill ranges are identical, so they share the information
carrier parameters, memory field sizes, and the rule for locating those
fields. It suffices to store, for each class, its size, the first memory
location, both child spill sizes, and the constants used by the lemma.

There are only constantly many classes per level, and the tree height is
\[
    O(\log N)=O(\log n).
\]
The total number of precomputed word constants is therefore \(O(\log n)\).
For heap index \(i\), its level is
\[
    \lfloor\log_2 i\rfloor,
\]
which is computable by a constant number of Word RAM operations. The level
then identifies the node's class and the shared parameters and memory
locations in constant time.

\subsection{The Final-Bit Issue}

The direct construction yields
\[
    \left\lceil n\log_2\Sigma+o(1)\right\rceil
\]
bits, which can in principle equal
\(\lceil n\log_2\Sigma\rceil+1\). The paper removes this possible last bit by
increasing the block, and hence carrier, universe. Arithmetic then uses
\(O(c\log n)\)-bit integers for a fixed constant \(c\), still a constant
number of Word RAM words, and the accumulated redundancy becomes
\(O(n^{-c})\).

For fixed \(\Sigma\), the paper invokes a lower bound from the theory of
linear forms in logarithms: unless \(n\log_2\Sigma\) is already integral, its
distance from an integer is at least inverse polynomial in \(n\). Choosing
\(c\) sufficiently large makes the remaining redundancy smaller than this
separation, ensuring
\[
    n\log_2\Sigma+O(n^{-c})
    \leq
    \left\lceil n\log_2\Sigma\right\rceil.
\]
Only this separation consequence is needed; the underlying number-theoretic
theorem is not reproved here.

\subsection{The Vector Representation Theorem}

\begin{theorem}
\label{thm:vector-representation}
On the Word RAM, a vector \(A[1\ldots n]\) over an alphabet of size
\(\Sigma\) can be represented using
\[
    \left\lceil n\log_2\Sigma\right\rceil
\]
bits while supporting reads and writes in \(O(1)\) time. The representation
requires \(O(\log n)\) precomputed word constants.
\end{theorem}

\begin{proof}
Grouping symbols into carriers of universe \(X=\Theta(n^2)\) reduces the sum
of local rounding redundancies to \(o(1)\). The tree preserves the lemma's
local recovery asymmetry, so a carrier is read from its node and parent and
an update is absorbed after re-encoding at most those two nodes. Nodes of the
same level and subtree class share parameters, leaving only \(O(\log n)\)
precomputed word constants. Finally, enlarging the carriers reduces the
residual error below the inverse-polynomial separation from the next integer,
which gives the exact \(\lceil n\log_2\Sigma\rceil\)-bit bound.
\end{proof}
